\subsection{Quantum spin chains and the lattice Lax equation}
\label{sec:quantum-lattice-setup}

The two lattice models studied in this paper are homogeneous periodic spin
chains with nearest-neighbour interactions and a two-dimensional local Hilbert
space.  We use the standard quantum inverse-scattering framework
\cite{Faddeev1996,AvanDoikouSfetsos2010}.  Let
$\mathcal H_n\simeq\mathbb C^2$ be the Hilbert space at site $n$ and
$\mathcal H^{(N)}=\bigotimes_{n=1}^{N}\mathcal H_n$ the Hilbert space of the
periodic chain.  A fixed two-site operator
$h\in\End(\mathbb C^2\otimes\mathbb C^2)$ is embedded as $h_{n,n+1}$ on
each nearest-neighbour bond, with site labels understood modulo $N$.  The
lattice Hamiltonian is
\begin{equation}
H^{\mathrm{lat}}
 = \sum_{n=1}^{N} h_{n,n+1},
\qquad h_{N,N+1}\equiv h_{N,1}.
\label{eq:generic-lattice-hamiltonian}
\end{equation}

For the integrable chains considered here, the local interaction is encoded by
a quantum $R$-matrix of difference form.  Let
$\mathcal V_a\simeq\mathbb C^2$ denote an auxiliary space and
$u\in\mathbb C$ the quantum spectral parameter.  We write $R_{ab}(u)$ and
$R_{a,n}(u)$ for the same two-site $R$-matrix acting on the indicated tensor
factors.  The $R$-matrix satisfies the Yang--Baxter equation and is normalized
at the regular point by
\begin{equation}
R(0)=P,
\qquad\text{equivalently}\qquad
R_{ab}(0)=P_{ab},
\label{eq:R-regularity}
\end{equation}
where
\begin{equation}
P\in\End(\mathbb C^2\otimes\mathbb C^2)
\label{eq:local-permutation-operator}
\end{equation}
is the permutation operator,
$P(|v\rangle\otimes|w\rangle)=|w\rangle\otimes|v\rangle$.
For a regular fundamental spin-$\tfrac12$ chain, the two-site Hamiltonian is
\begin{equation}
h
=
P\left.\frac{\partial R(u)}{\partial u}\right|_{u=0}.
\label{eq:h-from-R}
\end{equation}
This is the standard relation between a regular $R$-matrix and the
nearest-neighbour Hamiltonian density
\cite{Faddeev1996,deLeeuwPribytokRyan2019}.

A scalar rescaling $R(u)\mapsto f(u)R(u)$ with $f(0)=1$ preserves
Eq.~\eqref{eq:R-regularity} and adds $f'(0)\Id$ to $h$.  This shifts the
total energy by a constant and leaves the Heisenberg equations unchanged.  An
overall rescaling of the spectral parameter corresponds to a normalization of
the Hamiltonian.  These freedoms are fixed together with each explicit
$R$-matrix.

The quantum lattice Lax operator at site $n$ is
\begin{equation}
L_{a,n}(u):=R_{a,n}(u)
\in\End(\mathcal V_a\otimes\mathcal H_n).
\label{eq:lattice-Lax-operator}
\end{equation}
Differentiating the Yang--Baxter equation at the regular point and using
Eq.~\eqref{eq:h-from-R} gives the Sutherland relation
\cite{deLeeuwPribytokRyan2019},
\begin{align}
&[h_{n-1,n},L_{a,n}(u)L_{a,n-1}(u)]
\nonumber\\
&\qquad=
L_{a,n}(u)\,\partial_u L_{a,n-1}(u)
-
\bigl(\partial_uL_{a,n}(u)\bigr)L_{a,n-1}(u).
\label{eq:Sutherland-relation}
\end{align}
Here $h_{n-1,n}$ is the same two-site Hamiltonian acting on the neighbouring
physical sites $n-1$ and $n$.

Let $t_{\rm lat}$ denote the quantum lattice time.  For an operator
$\mathcal O$ with no explicit time dependence we use the Heisenberg equation
\begin{equation}
\partial_{t_{\rm lat}}\mathcal O
=\ii[H^{\rm lat},\mathcal O].
\label{eq:Heisenberg-lattice-time}
\end{equation}
Equation~\eqref{eq:Heisenberg-lattice-time} defines the algebraic time
evolution used throughout the lattice Lax calculation.  For a non-Hermitian
Hamiltonian, normalized expectation values of right Schr\"odinger states in
the standard inner product define a separate evolution problem.  The
continuum comparison below concerns the algebraic evolution in
Eq.~\eqref{eq:Heisenberg-lattice-time} and its coherent-state lower symbols.
Because $L_{a,n}$ acts only on the physical site $n$, only the two Hamiltonian
terms $h_{n-1,n}$ and $h_{n,n+1}$ contribute to its commutator with
$H^{\rm lat}$.  The Sutherland relation organizes these two contributions
into a discrete zero-curvature equation.

For spectral parameters at which $L_{a,n}(u)$ is invertible, define the
\emph{finite-lattice time Lax operator}
\begin{align}
A_{a,n}(u)
={}&\ii h_{n-1,n}
-\ii L_{a,n}(u)^{-1}h_{n-1,n}L_{a,n}(u)
\nonumber\\
&-\ii L_{a,n}(u)^{-1}\partial_uL_{a,n}(u).
\label{eq:finite-lattice-time-Lax}
\end{align}
The physical part of $A_{a,n}$ acts on the neighbouring sites
$n-1$ and $n$, while its matrix indices belong to the auxiliary space
$\mathcal V_a$.  Equations~\eqref{eq:Sutherland-relation} and
\eqref{eq:Heisenberg-lattice-time} then give
\begin{equation}
\partial_{t_{\rm lat}}L_{a,n}(u)
=
A_{a,n+1}(u)L_{a,n}(u)
-
L_{a,n}(u)A_{a,n}(u).
\label{eq:discrete-zero-curvature}
\end{equation}
This is the lattice analogue of a zero-curvature equation.  The time Lax
operator has a scalar ambiguity: adding to every $A_{a,n}$ the same
scalar function of $u$ times the identity on the full tensor product leaves
\eqref{eq:discrete-zero-curvature} unchanged.  This freedom permits convenient
normalizations.

The continuum limit of \eqref{eq:discrete-zero-curvature} depends on the
physical support of the operators in its two products.  They contain
operators that share a physical lattice site: $A_{a,n+1}$ and
$L_{a,n}$ both act on $\mathcal H_n$, while $L_{a,n}$ and
$A_{a,n}$ also overlap on $\mathcal H_n$.  Consequently, coherent-state expectation values of these products retain
correlations associated with the shared physical site.
